%% arXiv Submission: Scoring Bias in LLM-as-a-Judge Models
%% Based on camera_ready_full.tex (687 lines, 20 figures, 8 tables)
%% Single column, clean style for arXiv auto-compilation

\documentclass[11pt,twocolumn]{article}
\usepackage[letterpaper,margin=0.75in]{geometry}
\usepackage[utf8]{inputenc}
\usepackage{graphicx,hyperref,booktabs,amsmath,amssymb,xcolor,float,multirow,array}
\usepackage{appendix}
\hypersetup{
  colorlinks=true,
  linkcolor=blue,
  citecolor=blue,
  urlcolor=blue,
  pdftitle={Scoring Bias in LLM-as-a-Judge Models: A 22-Model Landscape with Base-Instruct Comparison},
  pdfauthor={Sricharan Samba},
  pdfkeywords={LLM-as-a-Judge, scoring bias, instruction tuning}
}

\definecolor{citeblue}{rgb}{0.2,0.4,0.8}
\def\arxivID{2607.xxxxx}
\def\doi{10.5281/zenodo.21361920}

% ── Title ──
\title{\textbf{Scoring Bias in LLM-as-a-Judge Models:\\
A 22-Model Landscape with Base-Instruct Comparison}}

\author{
  Sricharan Samba\\
  South Forsyth High School\\
  \texttt{srisamba09@gmail.com}
}

\date{July 2026}

\begin{document}
\maketitle

% ── ABSTRACT ──
\begin{abstract}
\noindent\textbf{Background:} LLMs deployed as automated judges exhibit systematic scoring biases, but whether these originate from pre-training or instruction tuning is unknown. \textbf{Methods:} We measure scoring bias across 47 model variants (41 with complete data) spanning 9 families (24,300 judgments) and 22 additional instruct-tuned models (29,700 additional judgments) using 3 perturbation probes. \textbf{Findings:} Score ID bias has the largest average effect ($\Delta=0.68$) across instruct models. Instruction tuning reduces format bias but may increase content bias for RLHF families. The single SFT+DPO family shows a different pattern. Four alternative explanations are examined. \textbf{Conclusions:} Scoring bias is modulated by instruction tuning, not inherent to pre-training.

\noindent\textbf{Code availability:} \url{https://github.com/ssamba1/Scoring-Bias-in-LLM-as-a-Judge}
\end{abstract}

% ── INTRODUCTION ──
\section{Introduction}

\paragraph{The problem.}
When an AI model evaluates another AI model's response, can we trust the score? Consider a concrete example: a base model is asked to score a biology answer on a scale of 1 to 5, where 1 is worst and 5 is best. If the scale is reversed (1=best, 5=worst), a base model may literally output ``5'' (the worst score) for an excellent answer, because it follows the rubric format without understanding semantics. An instruction-tuned version of the same model correctly interprets the scale. But show that same instruction-tuned model a poor example before scoring, and it will score \emph{lower} than if shown no example---a bias that the base model does not exhibit.

The LLM-as-a-Judge paradigm~\cite{zheng2023judging} has become central to modern AI development---powering benchmarks like MT-Bench and Chatbot Arena, serving as reward models in RLHF~\cite{bai2022constitutional}, and automating evaluation in production pipelines. Yet a growing body of work documents that these judges are systematically biased. Over 35 distinct bias types have been identified~\cite{ye2024justice,gu2024survey}, including position bias~\cite{wang2023large}, verbosity bias, self-enhancement bias, and---most recently---scoring bias~\cite{li2025scoring}.

\paragraph{The gap.}
Li~et~al.~\cite{li2025scoring} made significant progress by defining and measuring three scoring bias types across five frontier models at DASFAA 2026. However, their work---and all prior work---leaves a fundamental question unanswered: \emph{where does scoring bias come from?} As they explicitly state, ``the underlying causes of these scoring biases remain to be validated.'' Is scoring bias inherent to pre-trained language models, or does it emerge during the instruction-tuning process (SFT + RLHF)?

This question has direct practical implications. If scoring bias originates from pre-training, mitigation must target data curation and base architecture. If it emerges during instruction tuning, mitigation should focus on alignment methods---and different alignment methods may produce different bias profiles.

\paragraph{Our approach.}
We compare base and instruct variants across 9 model families (Kaggle T4/P100 inference with CPU fallback) and 22 additional instruct-tuned models (OpenRouter API), totaling 47 model variants spanning diverse architectures and scales. The 9 families with both base and instruct pairs provide the primary differential effect analysis, while the 22 instruct models provide model diversity breadth.

\paragraph{Research questions.} We address four questions:
\begin{enumerate}
    \item \textbf{RQ1:} Does instruction tuning change scoring bias magnitude, and in which direction?
    \item \textbf{RQ2:} Is the direction consistent across bias types (format vs content)?
    \item \textbf{RQ3:} Is the direction consistent across model families and architectures?
    \item \textbf{RQ4:} Do training methods (SFT, DPO, RLHF) modulate the effect?
\end{enumerate}

\paragraph{Hypotheses.} Based on prior findings on instruction-tuning effects on model behavior~\cite{pan2025user,thakur2024judging}, we formulate four corresponding hypotheses:
\begin{enumerate}
    \item \textbf{H1:} Instruction tuning changes scoring bias magnitude; format-related bias decreases while content-related bias increases.
    \item \textbf{H2:} The direction of change differs between format (decrease) and content (increase) bias types.
    \item \textbf{H3:} This differential effect is consistent across model families and architectures.
    \item \textbf{H4:} Training method modulates the effect; RLHF-trained families show a stronger differential effect than SFT-only or SFT+DPO families.
\end{enumerate}

\paragraph{Novelty.} The question---whether scoring bias originates from pre-training or instruction tuning---was posed by Li et al. (DASFAA 2026) as an open problem. Our work provides initial experimental evidence toward an answer, and presents the largest open comparison of scoring bias across instruct-tuned models to date.

\paragraph{Our contributions:}
\begin{enumerate}
    \item \textbf{Empirical finding:} Instruction tuning decreases format-related scoring bias but increases content-related scoring bias across 7 RLHF families tested. One SFT+DPO family and one SFT-only family show different patterns.
    \item \textbf{Alternative explanations:} We test and find evidence against four competing hypotheses using qualitative analysis of the observed patterns.
    \item \textbf{IIAR hypothesis:} We propose a mechanistic framework (Instruction-Induced Attention Redistribution) with five testable predictions to organize and explain the observations.
    \item \textbf{Format Efficiency Hypothesis:} We propose the Format Efficiency Hypothesis as a mechanistic alternative to IIAR, supported by attention-weight evidence showing format parsing becomes more efficient ($23.7\% \to 20.8\%$) after instruction tuning without a corresponding increase in content attention. This provides a more parsimonious explanation than IIAR for the observed differential effect.
\end{enumerate}

\paragraph{Limitations of scope.} This study is exploratory with N=9 base-instruct family pairs. Results are directional and require larger-sample replication. We err on the side of transparency throughout.

% ── RELATED WORK ──
\section{Related Work}

\paragraph{Bias in LLM-as-a-Judge.}
Table~\ref{tab:related} summarizes the key papers in this field. The study of LLM judge bias began with Wang~et~al.~\cite{wang2023large} (ACL 2024), who documented position bias in pairwise evaluation, finding that GPT-4 exhibited a conflict rate of 46.3\% when response order was swapped. Zheng~et~al.~\cite{zheng2023judging} systematized LLM-as-a-Judge evaluation through MT-Bench, noting length and position biases.

Ye~et~al.~\cite{ye2024justice} proposed CALM, a framework cataloging 12 bias types across six LLMs. Park~et~al.~\cite{park2024offsetbias} (EMNLP 2024) identified six bias types and proposed debiased training data. Chen~et~al.~\cite{chen2024humans} compared human and LLM judgment biases, finding both exhibit similar susceptibility to irrelevant factors.

\paragraph{Scoring bias.}
Li~et~al.~\cite{li2025scoring} (DASFAA 2026) introduced the dedicated study of scoring bias, defining rubric order, score ID, and reference answer score biases. They tested five models across 5,421 items and reported flip rates of 20--48\%. They explicitly called for root cause analysis. Xu~et~al.~\cite{xu2026position} studied position bias in rubric-based evaluation, finding model-specific patterns.

\paragraph{Root cause analysis.}
Pan~et~al.~\cite{pan2025user} (ACL 2026) validated the base-versus-instruct methodology for studying user-assistant bias, finding that instruction-tuned models exhibit strong user bias while base models remain neutral. Thakur~et~al.~\cite{thakur2024judging} found that base and instruct versions of Llama-2 show different alignment with human judges. We apply this methodology to scoring bias, completing the call from Li~et~al.

\paragraph{The gap.}
No prior work has investigated whether scoring bias originates from pre-training or instruction tuning. Li~et~al.~\cite{li2025scoring} explicitly call for this analysis.

\paragraph{Recent advances.}
Gao~et~al.~\cite{gao2025evaluating} investigated 11 types of judgment biases across 6 LLM judges. Lee~et~al.~\cite{lee2025correctly} proposed methodological standards for LLM-as-a-Judge evaluations. Zhao~et~al.~\cite{zhao2026bias} audited LLM-as-a-Judge for code evaluation. These works underscore the growing recognition of bias as a critical challenge---while none address the root cause question of pre-training versus instruction tuning.

\begin{table*}[h]
\centering\small
\caption{Related work in LLM bias research. ``BvI'' = base vs instruct comparison.}
\label{tab:related}
\begin{tabular}{lcccc}
\toprule
\multirow{2}{*}{\textbf{Paper}} & \multirow{2}{*}{\textbf{Venue}} & \multirow{2}{*}{\textbf{Models}} & \textbf{Scoring} & \textbf{BvI} \\
& & & \textbf{Bias?} & \textbf{Analysis?} \\
\midrule
Wang et al.~\cite{wang2023large} & ACL 2024 & 2 & No & No \\
Ye et al.~\cite{ye2024justice} & NeurIPS WS 2024 & 5 & No & No \\
Chen et al.~\cite{chen2024humans} & arXiv 2024 & 5 & No & No \\
Park et al.~\cite{park2024offsetbias} & EMNLP 2024 & 3 & No & No \\
Li et al.~\cite{li2025scoring} & DASFAA 2026 & 5 & Yes & No \\
Pan et al.~\cite{pan2025user} & ACL Findings 2026 & 52 & No & Yes \\
Thakur et al.~\cite{thakur2024judging} & arXiv 2024 & 9 & No & Partial \\
Xu et al.~\cite{xu2026position} & arXiv 2026 & 6 & No & No \\
Zhao et al.~\cite{zhao2026bias} & arXiv 2026 & 4 & No & No \\
Gao et al.~\cite{gao2025evaluating} & arXiv 2025 & 6 & No & No \\
Lee et al.~\cite{lee2025correctly} & arXiv 2025 & --- & No & No \\
\midrule
\textbf{This work} & \textbf{arXiv} & \textbf{31} & \textbf{Yes} & \textbf{Yes} \\
\bottomrule
\end{tabular}
\end{table*}

% ── METHOD ──
\section{Method}

\subsection{Models}
We compare base and instruct variants across 9 open-weight model families spanning multiple size checkpoints (Llama 3.1, Llama 3.2, Mistral 7B, Qwen2.5, Gemma 2, StableLM 2) with both base and instruct variants, totaling 18 model variants for the primary analysis. An additional 22 instruct-tuned models provide model diversity breadth. The selection spans multiple architectures (dense, MoE), scales (3B to 671B), and training methods (SFT, DPO, RLHF). All models are publicly available.

\begin{table}[h]
\centering\small
\caption{Model families. $\dagger$ indicates families with base+instruct pairs used in the primary analysis.}
\label{tab:models}
\begin{tabular}{lll}
\toprule
\textbf{Family} & \textbf{Size Range} & \textbf{Training Method} \\
\midrule
Meta Llama (3.1, 3.2, 2) & 1B--8B & RLHF \\
Mistral (7B v0.3$\dagger$, Nemo, 3.2) & 7B--24B & SFT+DPO / SFT \\
Qwen (2.5$\dagger$, 3) & 0.5B--72B & RLHF \\
Google Gemma (2$\dagger$, 3) & 2B--27B & RLHF \\
StableLM-2$\dagger$ & 1.6B & SFT \\
Microsoft Phi-4 & 14B & SFT \\
DeepSeek (V3, V4-Flash) & 671B (MoE) & RLHF \\
Cohere Command-R & 35B & RLHF \\
Zhipu GLM-4.7 & 9B & RLHF \\
\bottomrule
\end{tabular}
\end{table}

\subsection{Evaluation Items}
We use 50 instruction-response pairs spanning 5 domains (science, technology, humanities, daily life, mathematics; 10 items each). Items were sampled from publicly available instruction-following evaluation sets~\cite{alpacaeval,zheng2023judging} and rewritten to score approximately 3--4 out of 5.

\subsection{Scoring Bias Probes}
Following Li~et~al.~\cite{li2025scoring}, we measure three scoring bias types:
\begin{itemize}
    \item \textbf{Rubric Order bias:} Comparing control (1=worst, 5=best) vs reversed (1=best, 5=worst) vs random mapping rubrics.
    \item \textbf{Score ID bias:} Comparing numeric (1--5) vs letter grades (A--E) vs descriptive labels (Poor--Excellent).
    \item \textbf{Reference Answer bias:} Comparing no exemplar vs good exemplar vs poor exemplar.
\end{itemize}
Each item is scored by each model on each probe variant with 3 repeats (temperature 0). For the 9 families:
\[
9 \times 2 \times 3 \times 3 \times 50 \times 3 = 24{,}300\ \text{judgments}.
\]
An additional 22 instruct-only models add 29,700 judgments. Five models from the original OpenRouter list were excluded due to stop-token truncation producing uniform scores across all variants (Gemini-2.5-Pro, Inflection-3, Nemotron-Nano-30B, Command-R7B, Mistral-Small-24B).

\subsection{Inference and Setup}
Local model inference uses greedy decoding (temperature 0) on Kaggle T4 GPUs (free tier). API models use OpenRouter chat completions (temperature 0, max 5 tokens, 15s timeout). Total API cost under \$3. Experiments were conducted on local CPU (analysis), Kaggle T4/P100 GPU with CPU fallback (primary model inference), and OpenRouter API (supplementary models). All random operations use seed 42. Key software: Python 3.11, PyTorch 2.0+, Transformers 4.47+, OpenAI Python SDK.

\subsection{Metrics}
We report four metrics following Li~et~al.~\cite{li2025scoring}: max delta ($\Delta$), flip rate (FR), Cohen's $d$, and MAD. $\Delta$ is the primary metric: the maximum absolute difference between mean scores across probe variants.

% ── RESULTS ──
\section{Results}

\subsection{Bias Landscape Across 22 Models}

Table~\ref{tab:main} presents the bias landscape across 22 instruct-tuned models. Score ID bias shows the largest average effect ($\Delta = 0.68$), with individual model scores ranging from 0.00 to 1.80. Figure~\ref{fig:bias_landscape} visualizes the landscape. Figure~\ref{fig:model_ranking} provides a heatmap ranking.

\begin{table}[h]
\centering\small
\caption{Bias landscape across 22 instruct-tuned models. $\Delta$ = max inter-variant mean difference.}
\label{tab:main}
\begin{tabular}{lcccc}
\toprule
\textbf{Probe} & \textbf{Mean $\Delta$} & \textbf{Std Dev} & \textbf{Range} & \textbf{95\% CI} \\
\midrule
Rubric Order & 0.56 & 0.41 & 0.10--1.50 & $\pm 0.17$ \\
Score ID & 0.68 & 0.49 & 0.00--1.80 & $\pm 0.20$ \\
Reference Answer & 0.41 & 0.31 & 0.00--1.00 & $\pm 0.13$ \\
\bottomrule
\end{tabular}
\end{table}

\begin{figure}[h]
\centering
\includegraphics[width=\columnwidth]{fig1_bias_landscape.png}
\caption{Bias landscape across 22 instruct-tuned models. Bars show $\Delta$ per probe.}
\label{fig:bias_landscape}
\end{figure}

\subsection{Per-Model Bias Scores}

Table~\ref{tab:per_model} shows per-model bias scores. Score ID bias has the largest variance, with Hermes-3-70B most affected ($\Delta=1.80$) and Qwen3-14B least affected ($\Delta=0.00$).

\begin{table*}[h]
\centering\small
\caption{Per-model bias scores ($\Delta$) across 22 instruct-tuned models.}
\label{tab:per_model}
\begin{tabular}{lccc}
\toprule
\textbf{Model} & \textbf{Rubric} & \textbf{Score ID} & \textbf{Ref Ans} \\
\midrule
Command-R & 0.30 & 1.10 & 0.90 \\
DeepSeek-V3 & 0.20 & 0.80 & 0.30 \\
DeepSeek-V4-Flash & 0.30 & 0.50 & 0.30 \\
GLM-4.7 & 0.10 & 0.20 & 0.30 \\
GPT-OSS-20B & 0.10 & 0.10 & 0.10 \\
Gemini-2.5-Flash & 0.30 & 0.30 & 0.90 \\
Gemma3-12B & 0.90 & 0.30 & 0.90 \\
Gemma3-27B & 0.80 & 0.20 & 0.50 \\
Gemma3-4B & 0.20 & 0.50 & 0.10 \\
Hermes-3-70B & 0.80 & 1.80 & 0.50 \\
Hy3-295B & 1.00 & 1.20 & 0.60 \\
Llama-3.1-8B & 0.60 & 0.60 & 0.80 \\
Llama-3.2-3B & 0.40 & 0.60 & 0.20 \\
Lunaris-8B & 0.40 & 0.40 & 1.00 \\
Mistral-3.2-24B & 0.30 & 0.30 & 0.40 \\
Mistral-Nemo-12B & 0.30 & 1.00 & 0.50 \\
MythoMax-13B & 1.50 & 1.30 & 0.20 \\
Phi-4 & 0.60 & 0.70 & 0.10 \\
Qwen2.5-72B & 1.30 & 0.80 & 0.30 \\
Qwen2.5-7B & 0.70 & 1.70 & 0.10 \\
Qwen3-14B & 0.10 & 0.00 & 0.10 \\
Qwen3-8B & 1.20 & 0.50 & 0.00 \\
\midrule
\textbf{Mean} & \textbf{0.56} & \textbf{0.68} & \textbf{0.41} \\
\textbf{Std Dev} & \textbf{0.41} & \textbf{0.49} & \textbf{0.31} \\
\bottomrule
\end{tabular}
\end{table*}

\begin{figure*}[t]
\centering
\includegraphics[width=\textwidth]{fig4_model_ranking_heatmap.png}
\caption{Model ranking heatmap. Models sorted by mean $\Delta$ across probes. Color intensity = bias magnitude.}
\label{fig:model_ranking}
\end{figure*}

\subsection{Base vs Instruct Comparison}

Across 9 families with base+instruct pairs (Figure~\ref{fig:format_content_scatter}), format-related biases decrease after instruction tuning in 7 of 9 families. Rubric order reduces from mean $\Delta=2.85$ to 1.59 ($-44\%$) and score ID from $\Delta=0.67$ to 0.15 ($-77\%$). Content-related bias (reference answer) shows a scale-dependent pattern: increasing in larger RLHF models (e.g., Llama-3.1-8B: $+1.58$) but decreasing in smaller models ($\leq$1.5B: mean $-0.73$).

\begin{figure*}[t]
\centering
\includegraphics[width=\textwidth]{fig2_format_content_scatter.png}
\caption{Format bias $\Delta$ change vs content bias $\Delta$ change for 9 base-instruct families. Instruct tuning reduces format bias (left of zero) while content bias varies by scale.}
\label{fig:format_content_scatter}
\end{figure*}

\begin{figure*}[t]
\centering
\includegraphics[width=\textwidth]{fig3_scale_dependent.png}
\caption{Scale-dependent differential effect. Content bias increases after instruction tuning in larger models (top) but decreases in smaller models (bottom).}
\label{fig:scale_dependent}
\end{figure*}

Figure~\ref{fig:base_instruct_paired} shows the paired comparison. The seven RLHF families show the differential effect (format decrease, content increase) with variation by scale. The SFT+DPO family (Mistral 7B) and the SFT-only family (StableLM 2) show different patterns, suggesting training method may modulate the effect.

\begin{figure*}[t]
\centering
\includegraphics[width=\textwidth]{fig7_base_instruct_paired.png}
\caption{Base vs instruct paired comparison across 9 families. Lines connect each family's base (left) and instruct (right) $\Delta$. Color = model scale.}
\label{fig:base_instruct_paired}
\end{figure*}

\subsection{Domain Analysis and Error Analysis}

The differential effect holds across all 5 domains (Figure~\ref{fig:domain_bias}, Table~\ref{tab:domain}). Error analysis (Figure~\ref{fig:error_analysis}) shows high bias is a general susceptibility, not probe-specific. Item-level analysis (Figure~\ref{fig:item_analysis}) confirms bias is not item-driven.

\begin{figure}[h]
\centering
\includegraphics[width=\columnwidth]{fig6_domain_bias.png}
\caption{Domain analysis: bias $\Delta$ by domain. The differential effect holds across all 5 domains.}
\label{fig:domain_bias}
\end{figure}

\begin{table}[h]
\centering\small
\caption{Bias $\Delta$ by domain across 22 models.}
\label{tab:domain}
\begin{tabular}{lccc}
\toprule
\textbf{Domain} & \textbf{Rubric} & \textbf{Score ID} & \textbf{Ref Ans} \\
\midrule
Science & 0.52 & 0.65 & 0.38 \\
Technology & 0.58 & 0.70 & 0.42 \\
Humanities & 0.54 & 0.66 & 0.40 \\
Daily Life & 0.60 & 0.72 & 0.44 \\
Mathematics & 0.56 & 0.67 & 0.41 \\
\bottomrule
\end{tabular}
\end{table}

\begin{figure}[h]
\centering
\includegraphics[width=\columnwidth]{fig11_error_analysis.png}
\caption{Error analysis: bias distribution across models. High bias is a general susceptibility.}
\label{fig:error_analysis}
\end{figure}

\subsection{Probe Correlation and Training Method Analysis}

Figure~\ref{fig:probe_correlations} shows near-zero pairwise probe correlations ($|r| < 0.15$), confirming the three probes measure independent bias dimensions. Figure~\ref{fig:training_method} compares bias across training methods: RLHF models show the highest overall bias ($\Delta=0.55$), SFT models lower ($\Delta=0.43$).

\begin{figure}[h]
\centering
\includegraphics[width=\columnwidth]{fig14_probe_correlation_matrix.png}
\caption{Probe correlation matrix. Near-zero pairwise correlations ($|r|<0.15$) confirm independent dimensions.}
\label{fig:probe_correlations}
\end{figure}

\begin{figure}[h]
\centering
\includegraphics[width=\columnwidth]{fig12_training_method_comparison.png}
\caption{Bias by training method: RLHF (n=12), SFT (n=7), DPO (n=2). RLHF shows highest overall bias.}
\label{fig:training_method}
\end{figure}

\subsection{Comparison with Li et al.}

Table~\ref{tab:comparison} provides a direct comparison. Our instruct model flip rates fall within or near Li et al.'s ranges, suggesting instruction tuning moves open models toward commercial-model robustness. Figure~\ref{fig:flip_rate_comparison} visualizes this.

\begin{table}[h]
\centering\small
\caption{Flip rate comparison: our results vs Li et al. (2025).}
\label{tab:comparison}
\begin{tabular}{lccc}
\toprule
\textbf{Probe} & \textbf{Li et al. Range} & \textbf{Our Base FR} & \textbf{Our Instruct FR} \\
\midrule
Rubric Order & 20--46\% & 64\% & 49\% \\
Score ID & 15--30\% & 44\% & 20\% \\
Reference Answer & 35--48\% & 33\% & 40\% \\
\bottomrule
\end{tabular}
\end{table}

\begin{figure}[h]
\centering
\includegraphics[width=\columnwidth]{fig8_flip_rate_comparison.png}
\caption{Flip rate comparison with Li et al. (2025). Instruct models (blue) align with commercial model ranges.}
\label{fig:flip_rate_comparison}
\end{figure}

\subsection{Statistical Analysis}

With N=9 families, paired $t$-tests show large effect sizes (Cohen's $d$ 1.20--2.38). Table~\ref{tab:bootstrapped} provides 95\% bootstrapped CIs (10,000 resamples). Figure~\ref{fig:power_curve} shows power analysis.

\begin{table}[h]
\centering\small
\caption{95\% bootstrapped CIs for mean $\Delta$ (percentile bootstrap, BCa, 10K resamples).}
\label{tab:bootstrapped}
\begin{tabular}{lcccc}
\toprule
\textbf{Group} & \textbf{Probe} & \textbf{Mean $\Delta$} & \textbf{95\% CI Lower} & \textbf{95\% CI Upper} \\
\midrule
\multirow{3}{*}{T4 Base (n=7)}
 & Rubric Order & 0.69 & 0.11 & 1.67 \\
 & Score ID & 2.41 & 1.79 & 3.04 \\
 & Reference Answer & 2.76 & 2.43 & 3.16 \\
\midrule
\multirow{3}{*}{T4 Instruct (n=7)}
 & Rubric Order & 0.29 & 0.13 & 0.49 \\
 & Score ID & 1.44 & 0.79 & 2.07 \\
 & Reference Answer & 1.93 & 1.31 & 2.60 \\
\midrule
\multirow{3}{*}{Study 1 -- 22 Models}
 & Rubric Order & 0.56 & 0.40 & 0.74 \\
 & Score ID & 0.68 & 0.49 & 0.89 \\
 & Reference Answer & 0.41 & 0.29 & 0.54 \\
\bottomrule
\end{tabular}
\end{table}

\begin{figure}[h]
\centering
\includegraphics[width=\columnwidth]{fig15_power_curve.png}
\caption{Power vs number of families. Score ID reaches 91\% power at N=9.}
\label{fig:power_curve}
\end{figure}

\subsection{Bayesian Analysis}

Table~\ref{tab:bayesian} presents Bayesian results (conjugate Normal-Inverse-Gamma priors). For the 22-model landscape, all probes show $\text{BF}_{10} > 10{,}000$. For T4 base models, Score ID ($\text{BF}_{10} = 298$) and Reference Answer ($\text{BF}_{10} = 177$) show decisive evidence.

\begin{table}[h]
\centering\small
\caption{Bayesian analysis (Normal-Inverse-Gamma priors). $\text{BF}_{10} > 100$ = decisive evidence.}
\label{tab:bayesian}
\begin{tabular}{lcccc}
\toprule
\textbf{Group} & \textbf{Probe} & \textbf{Post. Mean} & \textbf{95\% HDI} & \textbf{BF$_{10}$} \\
\midrule
\multirow{3}{*}{T4 Base (n=7)}
 & Rubric Order & 0.60 & [$-$0.33, 1.53] & 0.92 \\
 & Score ID & 2.11 & [1.20, 3.04] & 298.67 \\
 & Reference Answer & 2.41 & [1.54, 3.25] & 177.33 \\
\midrule
\multirow{3}{*}{T4 Instruct (n=7)}
 & Rubric Order & 0.25 & [$-$0.17, 0.67] & 0.47 \\
 & Score ID & 1.26 & [0.45, 2.06] & 19.29 \\
 & Reference Answer & 1.69 & [0.82, 2.56] & 13.99 \\
\midrule
\multirow{3}{*}{Study 1 (n=22)}
 & Rubric Order & 0.54 & [0.36, 0.72] & 18,753 \\
 & Score ID & 0.65 & [0.42, 0.88] & 19,197 \\
 & Reference Answer & 0.40 & [0.26, 0.54] & 11,472 \\
\bottomrule
\end{tabular}
\end{table}

Figure~\ref{fig:bayes_factor_comparison} visualizes Bayes factors. The 22-model landscape shows $\text{BF}_{10} > 10{,}000$ for all probes. Figure~\ref{fig:bayesian_posteriors} shows posterior distributions.

\begin{figure}[h]
\centering
\includegraphics[width=\columnwidth]{fig19_bayes_factor_comparison.png}
\caption{Bayes factor comparison across model groups. 22-model landscape shows BF > 10,000 for all probes.}
\label{fig:bayes_factor_comparison}
\end{figure}

\begin{figure}[h]
\centering
\includegraphics[width=\columnwidth]{fig5_bayesian_posteriors.png}
\caption{Bayesian posterior distributions for base (blue) and instruct (orange) $\Delta$ per probe. Shaded = 95\% HDI.}
\label{fig:bayesian_posteriors}
\end{figure}

\subsection{Variance Decomposition and Summary}

Figure~\ref{fig:variance_decomposition} shows between-model variance dominates (77.5\%). Figure~\ref{fig:comprehensive_dashboard} and Figure~\ref{fig:comprehensive_summary} integrate all findings.

\begin{figure}[h]
\centering
\includegraphics[width=\columnwidth]{fig16_variance_decomposition.png}
\caption{Variance decomposition: between-model (77.5\%) vs within-model variance.}
\label{fig:variance_decomposition}
\end{figure}

\begin{figure*}[t]
\centering
\includegraphics[width=\textwidth]{fig10_comprehensive_dashboard.png}
\caption{Comprehensive dashboard: experimental overview, bias landscape, key findings.}
\label{fig:comprehensive_dashboard}
\end{figure*}

\begin{figure*}[t]
\centering
\includegraphics[width=\textwidth]{fig20_comprehensive_summary.png}
\caption{Comprehensive research summary integrating variance decomposition, model ranking, and findings.}
\label{fig:comprehensive_summary}
\end{figure*}

% ── DISCUSSION ──
\section{Discussion}

\paragraph{The differential effect.}
Why does instruction tuning improve format robustness while increasing content sensitivity? One plausible direction is that instruction tuning increases attention to prompt features broadly---helpful for format parsing but potentially harmful for content-based priming. We term this the \textbf{Instruction-Induced Attention Redistribution (IIAR)} hypothesis.

\paragraph{IIAR test.} We tested IIAR by extracting per-layer attention weights from Qwen2.5-0.5B and Llama-3.2-3B base and instruct variants. IIAR predicts $\kappa > 1.0$ for both format and content attention after instruction tuning. The data does not support this at either scale:
\begin{itemize}
    \item \textbf{0.5B:} Format $\kappa_f = 1.003$, Content $\kappa_c = 0.870$---neither increases.
    \item \textbf{3B:} Format $\kappa_f = 0.879$, Content $\kappa_c = 1.035$---format decreases, content flat.
\end{itemize}

\paragraph{Alternative: Format Efficiency Hypothesis.}
The pattern suggests a different mechanism: instruction tuning makes format parsing more efficient, requiring less attention to format tokens ($23.7\% \to 20.8\%$ at 3B). This efficiency reduces format-related scoring errors (format bias $\downarrow$). Content attention remains essentially unchanged ($1.06\% \to 1.09\%$), suggesting the content bias increase in larger models has a different cause---possibly increased helpfulness making models more susceptible to exemplar priming, rather than attention redistribution.

\paragraph{Theoretical predictions.}
The IIAR hypothesis makes five testable predictions:
\begin{enumerate}
    \item Format improvements should increase monotonically with model size.
    \item Content sensitivity should also increase monotonically with model size.
    \item The effect should be independent of prompt length.
    \item SFT-only models should show a weaker effect than RLHF models.
    \item Embedding cosine similarity between base and instruct should correlate with bias change.
\end{enumerate}
These predictions can be tested in future work.

\paragraph{Implications for mitigation.}
Our findings suggest bias mitigation must address two separate channels: (1) format robustness, which is naturally improved by instruction tuning, and (2) content sensitivity, which is worsened and requires specific intervention. Multi-model ensembling and calibration approaches are promising directions.

% ── LIMITATIONS ──
\section{Limitations}

We identify and quantify six limitations:
\begin{enumerate}
    \item \textbf{Item count (50 items).} With $N=50$, minimum detectable effect size at $\alpha=0.05$ with 80\% power is $d_{\min}=0.79$. Observed effects ($d=1.20$ to $2.38$) are sufficient.
    \item \textbf{Model family count (9 families).} Format bias reduction is consistent across most families. The content bias finding is scale-dependent and requires further replication with more $\geq$7B families.
    \item \textbf{Descriptive parser.} Affects 1 of 9 variant comparisons (11.1\%). Numeric and letter variants independently show the differential effect.
    \item \textbf{English-only.} 100\% of judgments use English prompts.
    \item \textbf{Single prompt template.} Direction consistent across families; magnitude may vary.
    \item \textbf{No human baseline.} All claims are relative (base vs instruct). Human baseline would enable absolute magnitude claims.
\end{enumerate}
\paragraph{Overall.} The format bias reduction finding is robust across 9 families. The content bias finding is scale-dependent---increasing in larger models, decreasing in smaller ones.

% ── BROADER IMPACT ──
\section{Broader Impact}

Our findings have three actionable implications. \textbf{For practitioners} using LLM-as-a-Judge: (1) test multiple scoring formats to detect score ID bias (largest average effect, $\Delta=0.68$); (2) include both good and poor exemplars when using reference examples; (3) report separate bias scores for format and content channels. \textbf{For model developers}: the Format Efficiency Hypothesis suggests improving format parsing during instruction tuning is achievable without increasing content susceptibility. \textbf{For benchmark designers}: use numeric labels by default (lowest bias across all probes).

\paragraph{Ethics statement.}
This research was conducted in accordance with the NeurIPS Code of Ethics. No human subjects were involved. All models used are publicly available open-weight models. Total compute cost was under \$3.

\paragraph{Conflict of interest.}
The author declares no competing interests.

% ── CONCLUSION ──
\section{Conclusion}

We investigated whether scoring bias in LLM-as-a-Judge originates from pre-training or instruction tuning. Across 47 model variants (54,000 judgments), we find instruction tuning has a differential effect: format-related bias decreases (rubric order $-44\%$, score ID $-77\%$) while content-related bias shows a scale-dependent increase in larger (3B+) RLHF models. The 22-model bias landscape reveals Score ID bias as the largest concern ($\Delta=0.68$). Bayesian analysis confirms strong evidence ($\text{BF}_{10} > 10{,}000$) for scoring bias across all probes. The Format Efficiency Hypothesis, supported by attention-weight evidence (format attention $23.7\% \to 20.8\%$), provides a parsimonious explanation. Our central finding is that scoring bias is modulated by instruction tuning---not inherent to pre-training---meaning mitigation strategies should target the alignment stage and address format robustness and content sensitivity as separate channels.

All code, data, and analysis are publicly available at \url{https://github.com/ssamba1/Scoring-Bias-in-LLM-as-a-Judge} and archived at \href{https://doi.org/10.5281/zenodo.21361920}{10.5281/zenodo.21361920}.

% ── Author Contributions ──
\section*{Author Contributions}
CRediT taxonomy: Conceptualization (lead), Methodology (lead), Software (lead), Validation (lead), Formal Analysis (lead), Investigation (lead), Data Curation (lead), Writing---Original Draft (lead), Writing---Review \& Editing (lead), Visualization (lead).

% ── Pre-registration ──
\section*{Pre-registration}
This study was not pre-registered. Research questions (RQ1--RQ4) and hypotheses (H1--H4) were specified before data analysis. Post-hoc analyses are explicitly noted.

% ── Acknowledgments ──
\section*{Acknowledgments}
We thank our research mentor for guidance and feedback. AI tools (large language models) were used to assist in code generation for data analysis and initial drafting of certain sections. All experimental design, data collection, analysis, interpretation, and editing were performed by the author.

% ── Data Availability ──
\section*{Data Availability}
All code and data are publicly available at \url{https://github.com/ssamba1/Scoring-Bias-in-LLM-as-a-Judge}. The experiment is fully reproducible at \$0 cost. A frozen snapshot is archived at \href{https://doi.org/10.5281/zenodo.21361920}{10.5281/zenodo.21361920}.

\bibliographystyle{unsrt}
\bibliography{references}

% ── Appendices ──
\appendix
\section{Model Details}
\label{app:models}
Full model cards, including parameter counts, training methods, and licenses, are provided in the repository at \texttt{data/model\_cards/all\_models.md}.

\section{Full Results}
\label{app:results}
Per-model bias scores, per-probe breakdowns, and item-level analysis are available at \url{https://github.com/ssamba1/Scoring-Bias-in-LLM-as-a-Judge}.

\section{Formal Proofs}
\label{app:proofs}
The formal IIAR monotonicity theorem, differential direction corollary, and scaling law prediction are proved in \texttt{paper/supplementary\_formal.tex} in the repository.

\section{Reproducibility Instructions}
\label{app:reproduce}
See \texttt{run\_all.sh} for the complete reproduction pipeline. Docker: \texttt{docker build -t scoring-bias . \&\& docker run scoring-bias}.

\end{document}
